\section{Method}
\label{sec:method}

\subsection{Problem Setup}

Let $x$ be a test image and let a frozen CLIP/AnomalyCLIP encoder produce
patch features $\{f_i\}_{i=1}^{N}$, where $f_i\in\mathbb{R}^{C}$. The normal
and abnormal text prototypes are denoted by $t_n$ and $t_a$. All features are
normalized before similarity computation. The initial local anomaly probability
is
\begin{equation}
  s_i^0 =
  \frac{\exp(\langle \bar f_i,\bar t_a\rangle/\tau)}
       {\exp(\langle \bar f_i,\bar t_n\rangle/\tau)+
        \exp(\langle \bar f_i,\bar t_a\rangle/\tau)} .
  \label{eq:initial-score}
\end{equation}
This is the same normal/abnormal prototype scoring view used by CLIP-based
ZSAD methods, but here $s_i^0$ is treated only as a semantic prior.

\subsection{Boundary-Aware Local Reliability}

To obtain a local signal that is not identical to the text-prototype response,
we compute a Haar-style decomposition on the CLIP patch-feature grid. Wavelet
multi-resolution analysis is a classical signal representation tool~
\cite{daubechies1988orthonormal,mallat1989theory}; we only use the simplest
feature-grid decomposition as a reliability definition. For each $2\times2$
neighborhood $(F_{00},F_{01},F_{10},F_{11})$ in the patch grid,
\begin{align}
  LL &= \frac{1}{2}(F_{00}+F_{01}+F_{10}+F_{11}), \\
  LH &= \frac{1}{2}(F_{00}-F_{01}+F_{10}-F_{11}), \\
  HL &= \frac{1}{2}(F_{00}+F_{01}-F_{10}-F_{11}), \\
  HH &= \frac{1}{2}(F_{00}-F_{01}-F_{10}+F_{11}).
\end{align}
The high-frequency magnitude is
\begin{equation}
  HF_i=\operatorname{mean}_c(|LH_{i,c}|+|HL_{i,c}|+|HH_{i,c}|).
\end{equation}
The low-frequency component provides a structure boundary estimate,
\begin{equation}
  E_i=\|\nabla LL_i\|_2 .
\end{equation}
After upsampling to the original patch grid, both maps are clipped by per-image
1/99 percentiles and normalized to $[0,1]$, giving $\widehat{HF}_i$ and
$\widehat{E}_i$. The reliability signal is
\begin{equation}
  W_i=\operatorname{Norm}\left(\widehat{HF}_i(1-\widehat{E}_i)\right).
  \label{eq:reliability}
\end{equation}
$W_i$ is not an anomaly score. It downweights high-frequency responses that
coincide with low-frequency structure boundaries, and is used only to guide
which patches become calibration evidence.

\subsection{Patch Evidence Selection}

We construct abnormal-side and normal-side evidence weights as
\begin{align}
  \omega_i^a &=
  (s_i^0)^\gamma\left((1-\lambda)+\lambda W_i\right)^\eta, \\
  \omega_i^n &=
  (1-s_i^0)^\gamma\left((1-\lambda)+\lambda(1-W_i)\right)^\eta .
  \label{eq:evidence-weights}
\end{align}
Thus abnormal evidence must be supported by the semantic abnormal prior and the
local reliability signal, while normal evidence must have low abnormal prior
and low abnormal-style reliability. For each image we select the top-ratio
patches under each weight and compute weighted visual prototypes:
\begin{align}
  v_a &= \operatorname{Norm}\left(
  \sum_{i\in\mathcal{A}_x}
  \frac{\omega_i^a}{\sum_{j\in\mathcal{A}_x}\omega_j^a} f_i\right), \\
  v_n &= \operatorname{Norm}\left(
  \sum_{i\in\mathcal{N}_x}
  \frac{\omega_i^n}{\sum_{j\in\mathcal{N}_x}\omega_j^n} f_i\right).
  \label{eq:visual-prototypes}
\end{align}
The sets $\mathcal{A}_x$ and $\mathcal{N}_x$ are not labels. They are
image-internal evidence sets selected by the current scoring rule.

\subsection{Conservative Prototype Calibration}

The visual prototypes are used to calibrate the text prototypes:
\begin{align}
  t_n^x &= \operatorname{Norm}\left((1-\beta_x)t_n+\beta_x v_n\right),\\
  \hat t_a^x &= \operatorname{Norm}\left((1-\alpha_x)t_a+\alpha_x v_a\right).
\end{align}
Because abnormal evidence is less reliable in a single test image, the update
is conservative. In the implementation, $\beta_x$ is scaled by the normal
evidence confidence and is applied only when abnormal evidence passes a
minimum-confidence gate; $\alpha_x$ is additionally gated by an abnormal
confidence threshold. The stable controlled setting uses
\[
\gamma=1,\quad \eta=1,\quad \lambda=0.05,\quad
\text{top-ratio}=0.20,
\]
\[
\alpha_0=0,\quad \beta_0=0.01,\quad \tau_a=0.15,\quad
\delta_a=0.06 .
\]
With $\alpha_0=0$, the abnormal prototype is effectively not updated in the
current setting. The method mainly estimates a small image-conditioned normal
reference and avoids drifting toward unreliable abnormal patches.

The final patch anomaly score is recomputed with the calibrated prototypes:
\begin{equation}
  s_i^x =
  \frac{\exp(\langle \bar f_i,\bar t_a^x\rangle/\tau)}
       {\exp(\langle \bar f_i,\bar t_n^x\rangle/\tau)+
        \exp(\langle \bar f_i,\bar t_a^x\rangle/\tau)} .
  \label{eq:final-score}
\end{equation}
When multiple CLIP patch layers are used, each layer is rescored with the same
image-conditioned prototypes and the layer maps are fused. The added module
uses no optimizer step, no gradient backpropagation, and no target-category
normal or abnormal labels.

\subsection{Method Boundary}

The method is intentionally different from direct local-cue fusion. A direct
fusion baseline forms an anomaly map by mixing $s_i^0$ and $W_i$ at the map
level. That baseline tests whether the local reliability map itself defines
anomaly. \method{} instead uses $W_i$ only before scoring, during evidence
selection and prototype calibration, and the final map remains a CLIP semantic
similarity map.
